% 本文件由 LaTeX 排版 Agent 根据有效 translation.json 自动生成。
% author=Rufinus; work_id=2711; title=宗徒信经注释

\chapter{宗徒信经注释}

% structure: p0001 source=h1 target=omitted reason=page-title-already-rendered-by-parent

本信经注释系应\Person{劳伦斯}主教之请而作，其教座所在不详；但\Person{丰塔尼尼}氏于所著《\Person{鲁菲努斯}传》中推测为\Place{康科迪亚}，即\Person{鲁菲努斯}之故乡。\par

其确切写作年代无法确定；但鉴于鲁菲努斯并未如在其若干著作中那样提及久居东方后书写拉丁文之困难，且文风相对流畅，故极可能写于其滞留\Place{阿奎莱亚}之晚年，即约三〇七至三〇九年间。\par

其价值相当可观：(一) 作为五世纪初地方教会中信经状态的见证，尤其各教会的变体。（在阿奎莱亚教会中，\OriginalTerm{在耶稣基督内}{in Jesu Christo}，\OriginalTerm{无形可见、不受苦难的圣父}{Patrem invisibilem et impassibilem}，\OriginalTerm{此肉身的复活}{Resurrectio hujus carnis}）；(二) 显示东方思想如何适应西方神学的形成；(三) 提供圣经正典的准则。\par

此注释清晰合理；除极少数段落（如从凤凰论证主的童贞诞生）外，至今仍对我们有用。\par

% structure: p0006 source=h2 target=section reason=relative-heading-rank; base=h2

\section{阿奎莱亚信经}

有必要提供\Person{鲁菲努斯}所注释的阿奎莱亚信经的拉丁文原文及英文译文。下列信经中特有的词或字母以斜体标出。\par

\begin{itemize}
\item \OriginalTerm{我信全能天主圣父，\Emph{无形可见，不受苦难}}{Credo in Deo Patre omnipotenti \Emph{invisibili et impassibili}}
\end{itemize}

\begin{itemize}
\item \OriginalTerm{又信耶稣基督，祂的独生子，我们的主}{Et in Jesu Christo, unico Filio ejus, Domino nostro}
\end{itemize}

\begin{itemize}
\item \OriginalTerm{祂由圣神降孕，生于童贞玛利亚}{Qui natus est de Spiritu Sancto ex Maria Virgine}
\end{itemize}

\begin{itemize}
\item \OriginalTerm{在本丢·彼拉多执政时被钉十字架，被埋葬}{Crucifixus sub Pontio Pilato, et sepultus}
\end{itemize}

\begin{itemize}
\item \OriginalTerm{祂下降地狱；第三日从死者中复活}{Descendit ad inferna; tertia die resurrexit a mortuis}
\end{itemize}

\begin{itemize}
\item \OriginalTerm{祂升天，坐在圣父之右}{Ascendit in cœlos; sedet ad dexteram Patris}
\end{itemize}

\begin{itemize}
\item \OriginalTerm{祂将从那里降来，审判生者死者}{Inde venturus est judicare vivos et mortuos}
\end{itemize}

\begin{itemize}
\item \OriginalTerm{又信圣神}{Et in Spiritu Sancto}
\end{itemize}

\begin{itemize}
\item \OriginalTerm{圣教会}{Sanctam Ecclesiam}
\end{itemize}

\begin{itemize}
\item \OriginalTerm{罪过的赦免}{Remissionem peccatorum}
\end{itemize}

\begin{itemize}
\item \OriginalTerm{\Emph{此}肉身的复活}{\Emph{Hujus} carnis resurrectionem}
\end{itemize}

1. 我心中既无写作之意，亦无写作之力，最忠信的主教（\OriginalTerm{教父}{Papa}）劳伦斯，因为我深知，以一己薄弱之能，置于公众批评之下，实非小可之险。但既然你在信中贸然（请恕我直言）以基督的圣事——我至为尊崇——要求我照信经的传统及自然意义，为你撰写有关信仰之事，虽然如此，你实加给我力不能胜的重担（因为我未忘哲人之言，其言极是：\Quote{论及天主，即使言说真理，亦属危险}）；然而，若你以祈祷助我因你要求而必须承担之责，我将试言一二——此举出于对你嘱托的敬畏，而非自信心有馀裕。然我所写恐难博得成熟读者之关注，更适于儿童及基督内初学者的理解。\par

我确知有些卓越的作者已就此问题出版过虔诚而简短的论文。此外，我知道异端者\Person{福提努斯}也写过同样主题；但其目的并非向读者解释经文的意义，而是将平实真实的话语曲解以支持自己的教义，然而圣神已确保在这些话语中未记载任何含糊、晦涩或与其他真理相悖之处：因为其中应验了那预言：\BibleQuote{要简约而彻底地在公义中成就言语：因为主要在地上作成简短的话。}我们当努力，首先恢复并强调宗徒话语本身的纯朴；其次，补足先前注释者似乎遗漏之处。但为使这所谓\Quote{简短的话}的范围更加明晰，我们将从起初探究它是如何赐予各教会的。\par

2. 我们的祖先传下传统：主升天之后，藉圣神降临，火焰似的舌头落在每位宗徒头上，使他们能说各种语言，以致无论多么遥远的民族、多么野蛮的言语，都不能对之封闭、不得接近；他们奉主命分别往各族去宣讲天主的圣言。因此，即将彼此分离之际，他们首先共同商定了将来宣讲的标准，以免分散后，在对那些被邀请信仰基督的人宣讲时，有任何不一致之处。于是全体聚集，充满圣神，如我们所说，他们编成了这份将来宣讲的简短纲要，每人各将自己的句子贡献给一个共同的总纲；并规定此规则应授予那些信从的人。\par

对此信式，基于许多且最充分的理由，他们称之为\OriginalTerm{符号}{\GreekText{κύμβολον}}。因为希腊文的\OriginalTerm{符号}{\GreekText{κύμβολον}}相当于拉丁文的\Quote{\OriginalTerm{标记}{Indicium}}（一个记号或标记）和\Quote{\OriginalTerm{共同贡献}{Collatio}}（由多人共同捐助）。因为宗徒们在这些话中，每人贡献了自己的句子。它被称为\Quote{\OriginalTerm{标记}{Indicium}}或\Quote{\OriginalTerm{标记}{Signum}}，一个记号或标记，因为，如宗徒保禄所说，以及\BibleBook{}{宗徒大事录}所载，那时许多四处游荡的犹太人，假装是基督的宗徒，为牟利或糊口而四处宣讲，虽提及基督之名，却不按精确的传统路线传递信息。因此宗徒们规定此信式作为一个记号或标记，使真正按宗徒规则宣讲基督的人可以被认出。最后，他们说，在内战中，由于双方盔甲相同，语言相同，战争的习惯和方式也相同，每位将军为防范背叛，通常向他的士兵交付一个独特的符号或口令——拉丁文称为\Quote{\OriginalTerm{标记}{Signum}}或\Quote{\OriginalTerm{标记}{Indicium}}——这样，如果遇到一个人，不确定他属于哪一边，询问符号（口令），他就会揭示他是友是敌。因此，传统继续，信经不写在纸或羊皮纸上，而是保留在信徒的心中，以确保没有人是通过阅读学会的，如同有时不信者那样，而是通过从宗徒传来的传统。\par

因此宗徒们，正如我们所说，即将分散去宣讲福音，确立了这一他们信仰一致的记号或标记；与诺厄的儿子们不同，他们即将彼此分离时，建造了一座用砖和沥青筑成的塔，塔顶可通天，而宗徒们竖起了一座信德纪念碑，由主的活石和珍珠组成，能抵抗敌人，既不被风吹倒，也不被洪水冲垮，更不被暴风骤雨的力量动摇。因此，前一代人在分离前夕建造了一座骄傲之塔，被罚口音混乱，以致无人能听懂邻人的言语；而后一代人建造了一座信德之塔，被赋予各种语言的知识和理解；这样，前者成为罪的标记，后者成为信德的标记。\par

但现在我们该对这些珍珠本身说些什么了，其中首要的是万有的泉源和源头，即所说的——\par

3. \Emph{我信全能者天主圣父。}\par

但在开始讨论这些话的意义之前，我认为最好提一下，在不同的教会中，这一条款中发现有一些添加。然而，在罗马城的教会中并非如此；原因，我猜想，一方面是因为那里没有异端起源，另一方面是因为那里保持着古老的习惯，即那些即将受洗的人应公开背诵信经，也就是说，在众人面前；其结果是，那些已是信徒的耳朵不会容忍添加一个字。但在其他地方，据我所知，由于某些异端者，似乎有了添加，希望通过这些添加排除教义上的新奇。然而，我们遵循我们在阿奎莱亚教会领洗时所接受的次序。\par

因此，\Emph{我信}放在首位，如宗徒保禄在\BibleBook{希伯来}{书}中所说：\BibleQuote{凡来到天主面前的，必须先相信祂存在，且相信祂是那些相信祂的人的赏报者。}先知也说：\BibleQuote{除非你们相信，你们不会明白。}因此，为使你们能理解，你们首先正确宣认你们相信；因为没有人会出海，将自己托付给深邃流动的元素，除非他首先相信航行可能安全；农夫也不会将种子投入犁沟，将谷粒撒在地上，除非相信雨水会来，连同太阳的温暖，在其滋养下，借着顺风，大地会生产、增多并使果实成熟。总之，生活中没有事情可以进行，除非首先有相信的意愿。那么，如果来到天主面前，我们首先宣认相信，这有何奇怪呢？因为没有相信，连普通生活也无法进行。我们在开头预先说了这些，因为外邦人常常反对我们说，我们的宗教因为缺乏理由，只依靠相信。因此，我们表明了，除非相信先行，否则任何事情既不可能做成，也不可能保持稳定。最后，婚姻是在相信会有孩子的情况下缔结的；孩子被托付给导师照料，是相信导师的教导会传递给学生；一人接受帝国的徽章，是相信人民、城市和装备精良的军队也会服从他。但如果在这些不同事业中，没有人着手进行，除非相信所述的结果会随之而来，那么，如果要认识天主，相信岂不更是必要的吗？但让我们看看信经中这个\Quote{简短话语}阐述了什么。\par

4. \Quote{\Emph{我信全能者天主圣父。}}\par

东方教会几乎普遍这样传授该条款：\Quote{我信\Emph{一个}全能者天主圣父；}而在下一条款中，我们说：\Quote{并且信基督耶稣，祂的独生子，我们的主，}他们则传授为：\Quote{并且信\Emph{一个}（主）我们的主耶稣基督，祂的独生子；}即宣认\Quote{一个天主}和\Quote{一个主}，依照宗徒保禄的权威。但我们稍后将回到这一点。目前，让我们注意这些话：\Quote{在全能者天主圣父内。}\par

\Quote{天主}，就人的心智所能形成的概念而言，是指那本性或实体，它超越万有之上。\Quote{圣父}是一个表达隐秘而不可言喻的奥秘的词。当你听到\Quote{天主}一词时，你应理解那是一种无始无终、单纯、不复合、不可见、无形体、不可言说、不可估量的实体，其中没有任何增添或受造之物。因为祂是万有的绝对原因，祂自身无因。当你听到\Quote{圣父}一词时，你应理解为某一位圣子的圣父，这圣子是前述实体的肖像。因为正如无人被称为\Quote{主}，除非他拥有财产或仆人，他是他们的主；正如无人被称为\Quote{师傅}，除非他有门徒；同样，无人可能被称为\Quote{父亲}，除非他有圣子。因此，\Quote{圣父}这个名号本身显明：圣父与圣子同存。\par

但我不愿你讨论天主圣父如何生了圣子，也不要过分好奇地窥探那深奥的奥秘，免得你过于急切地窥探不可接近的光明之辉煌，反而失去那藉天主恩宠赐予凡人的微弱一瞥。或者，如果你认为这是必须尽一切努力探究的题目，那么先向自己提出关乎我们自身的问题，然后，如果你能妥善处理它们，再从世俗之事飞升至天上之事，从可见的进至不可见的。首先，你若能，就确定：在你内的心智如何生成一个言语，以及其中记忆的性灵是什么；这些东西虽然在实在和运作上相异，但在本性或实质上又如何是一体；虽然它们出自心智，却从未与心智分离。如果这些东西虽然在我们之内、在我们自身灵魂的实质中，却因不能被肉眼看见而似乎向我们隐藏，那么让我们以更明显可见的事物作为探究对象。泉水如何从自身生出河流？藉着什么性灵它被带向湍急的溪流？如何发生：虽然河流与泉水是一体且不可分离，但河流既不能被理解为泉水，也不能被称为泉水，泉水也不是河流，然而看见河流的人也已看见了泉水？首先练习解释这些，你若能，就解释你手头的事物；然后你才能上升到更高的事理。然而，不要以为我要你一下子从地上升到天上：首先，若你允许，我要你注意这我们肉眼所见的穹苍，并请你解释，你若能，这可见之光的性质——那天火如何从自身生出光之辉煌，又如何产生热；虽然它们在实在上是三个，但在实质上如何是一。而如果你有能力探究其中的每一个，那么你仍须承认：神圣的生育奥秘是如此更加不同、更加超越，正如创造者比受造物更强大，工匠比他的作品更卓越，那永远存在者比那从无中开始存在者更崇高。\par

那么，天主是祂独生圣子我们的主之父，这是应相信的，而非应讨论的；因为仆人无权争论他主人的出生。圣父从天上作证说：\BibleQuote{这是我的爱圣子，我所喜悦的，你们要听祂。}圣父说祂是祂的圣子，并命令我们听祂。圣子说：\BibleQuote{看见我的，就是看见了圣父}，以及\BibleQuote{我与圣父原是一体}，以及\BibleQuote{我从天主出来而来到世上。}谁能自居于圣父与圣子这些话之间作为争论者？谁能分割神性，分离其意志，分裂其实质，将圣神切成部分，并否认真理所说的为真？因此，天主是真实的圣父，作为真理之父，不是从外部生育，而是从祂自身生出圣子；也就是说，作为全智者，祂生出智慧；作为公义者，生出公义；作为永远者，生出永远；作为不死者，生出不死；作为不可见者，生出不可见者；因为祂是光，生出辉煌；祂是心智，生出圣言。\par

5. 我们说过，东方教会在传授信经时说：\Quote{唯一的天主，全能的圣父}和\Quote{唯一的主}，这\Quote{一}不应从数量上理解，而是绝对地理解。例如，若有人说\Quote{一个人}或\Quote{一匹马}，这里的\Quote{一}是数量上的用法。因为可以有第二个人和第三个人，或第二匹马和第三匹马。但在不能添加第二或第三之处，若我们说\Quote{一}，我们意指的并非数量上的，而是绝对的一。例如，若我们说\Quote{一个太阳}，这里的意思是第二或第三不能被添加，因为只有一个太阳。更何况天主，当祂被称为\Quote{一}时，被称为\Quote{一}不是从数量上，而是绝对地，也就是说，祂之所以被称为一，是因为没有别的。同样，对主也应如此理解：祂是唯一的主，耶稣基督，天主圣父藉着祂或通过祂掌管万有，因此下一句称天主为\Quote{全能者}。\par

天主被称为\Emph{全能者}，因为祂掌管并统治万有。但圣父藉着圣子掌管万有，正如宗徒所说：\BibleQuote{藉着祂，万有被造，可见的与不可见的，无论是宝座、宰制者、首领或掌权者。}又于\BibleBook{}{希伯来书}中写道：\BibleQuote{藉着祂也创造了诸世界，}并\BibleQuote{立祂为万有的继承者。}所谓\Quote{立}我们应理解为\Quote{生}。如今，如果圣父藉着祂创造了诸世界，万有都是藉着祂被造，而祂又是万有的继承者，那么祂也藉着祂掌管万有。因为如同光生于光，真理生于真理，同样，全能者生于全能者。正如\BibleBook{若望}{默示录}中关于\Person{色辣芬}的记载：\BibleQuote{他们昼夜不得安息，不停地呼喊：圣、圣、圣、万军的上主，昔在、今在、将来永在的全能者。}那么，这位\BibleQuote{将来永在的}便被称为\BibleQuote{\Emph{全能者}}。除了基督、天主子，还有谁是\BibleQuote{将来永在的}呢？\par

在上述内容之后，加入了\Quote{\Emph{不可见与不可受苦}}。我应当指出，这两个词并不存在于罗马教会的信经中。众所周知，它们是为了对抗沙贝利异端而被加入我们教会的，此异端我们称之为\Quote{父受苦说}，即宣称圣父自己由童贞女所生而成为可见的，或主张祂在肉身内受苦。因此，为了排除对圣父的如此不虔敬，我们的先辈似乎加上了这些词，称圣父为\Quote{不可见与不可受苦}。因为显而易见，降生成人并在肉身内出生的是圣子，而非圣父；并且从肉身出生开始，圣子便成为\Quote{可见与可受苦的}。然而，就祂所拥有的、与圣父同一而不朽的天主性本体而言，我们必须相信圣父、圣子、圣神都不是\Quote{可见或可受苦的}。但圣子，因祂屈尊取了肉身，便在肉身内被看见并受苦。先知也预言了这一点，说道：\BibleQuote{这就是我们的天主，没有别的可比拟祂。祂寻得了知识的一切道路，将它赐给了祂的仆人雅各伯和祂所爱的以色列。此后，祂在地上显现，与人交谈。}\par

6. 随后接着是：\Quote{\Emph{又信基督耶稣，祂的独生子，我们的主。}}\Quote{耶稣}是希伯来语，意为\Quote{救主}。\Quote{基督}源于\Quote{傅油}，即受膏。因为在\BibleBook{}{梅瑟五书}中，我们读到\Person{纳威的儿子奥些}，当他被选为领导民众时，名字从\Quote{奥些}改为\Quote{耶稣}，以表明这是王侯和将帅的合宜名字，即那些应\Quote{拯救}跟随他们的人民的人。因此，两者都被称为\Quote{耶稣}：一位引导那从埃及地领出、从旷野漂泊中得释放的人民进入应许之地；另一位引导那从无知的黑暗中领出、从世界的谬误中召回的人民进入天国。\par

\Quote{基督}是大司祭或君王的专名。因为从前大司祭和君王都用傅油膏来祝圣；但他们作为可朽坏的人，用的是物质和可朽坏的膏。耶稣成为基督，是因被圣神所傅油，正如圣经论及祂所说：\BibleQuote{圣父以从天降下的圣神傅了祂。}依撒意亚也曾预表同样的事，以圣子的语气说：\BibleQuote{上主的神在我身上，因为祂给我傅了油，派遣我去向贫穷人传报喜讯。}\par

既然已指明\Quote{耶稣}是什么——祂拯救自己的子民，以及\Quote{基督}是什么——祂成为永远的大司祭，现在让我们看后续经文中这些话指谁：\Quote{祂的独生子，我们的主。}这里教导我们，我们所说的这位耶稣、我们所解释其名义的这位基督，是\Quote{天主的独生子}和\Quote{我们的主。}免得你或许认为这些人性的名称具有属世的意义，因此加上祂是\Quote{天主的独生子，我们的主。}因为祂是独一者所生的独一者，如同光有唯一的光辉，理智有独一的言语。非物质的生育并不落入复数，也不遭受分割，因为被生者与生者绝不分离。祂是\Quote{独一的}，如同思想对于心智，智慧对于智者，言语对于理智，勇毅对于勇士。因为正如圣父被宗徒称为\BibleQuote{独一智慧}，同样，独有圣子被称为智慧。因此，祂是\Quote{独生子。}并且，虽然在光荣、永恒、德性、主权、能力方面，祂与圣父相同，但这一切祂并非像圣父那样无始而拥有，而是出于圣父，作为圣子，无始而平等；虽然祂是万有的元首，但圣父是祂的元首。因为经上记载：\BibleQuote{基督的元首是天主。}\par

7. 当你听到\Quote{子}这个词时，不可想到肉身的出生；而要记住这是对非物质的、单纯而不复杂的本体的表述。因为，如我们上文所说，无论是理智生出言语，心智生出感觉，还是光从自身产生光明，在这一类生育中都不应寻求任何软弱或不完善的形象，那么我们对于这一切的创造者的观念，岂不应当更加纯净和神圣？\par

但或许你会说：\Quote{你所谈论的生成是一种无实体的生成。因为光不会产生有实体的光明，理智也不会产生有实体的言词；然而，据说天主圣子是实体性地生成的。} 对此我们首先答复：当在其他事物中使用比喻或例证时，相似之处不可能在每个细节上都成立，而只会在为使用该例证而设定的某一点上成立。例如，当在\WorkTitle{福音}中记载说：\BibleQuote{天国好比酵母，一个妇人把它藏在三斗面里}，难道我们认为天国在各方面都像酵母，以至于像酵母一样是可触摸、可腐朽、会变酸、不能使用的吗？显然，引用这个例证仅仅是为了说明：通过天主圣言的宣讲（这看似微小）如何能使人的心灵充满信德的酵母。同样，当记载说：\BibleQuote{天国好比一张撒在海里的网，网罗各种鱼}，难道我们假设天国的实体在各方面都类似于网线的质地以及网眼的结扣吗？不，比较的唯一目的是要表明：正如网将鱼从深海带到岸边，同样，通过天国福音的宣讲，人的灵魂从今世错误的深渊中被解放出来。由此可见，比喻或例证并不在被用来示范或说明的事物的每一个细节上都相符。否则，如果它们在各方面都相同，就不再被称为比喻或例证，而是事物本身了。\par

8. 此外还应注意到，没有一个受造物能与其创造者相同。因此，既然天主的本体或本质不容比较，神性也是如此。再者，每一个受造物都出自虚无。因此，如果一个火花（尽管如此无实体，但仍是火）从自身生出属于虚无的受造物，并在其中保持其源本质性的本性（即作为父亲火花之火），那么，那永恒之光（它始终存在，因为在其内没有非实体的东西）的本体为何不能从自身产生有实体的光明呢？因此，圣子被正确地称为\Quote{独一}的、\Quote{唯一}的。因为如此诞生的那位是\Quote{独一}且\Quote{唯一}的。唯一的事物不容比较。也不能说创造万物的那一位与其所创造的事物本质相似。因此，这就是耶稣基督，天主的独生子，也是我们的主。\Quote{独一}可以同时指圣子和主。因为耶稣基督作为真正的圣子和唯一的主，是\Quote{独一}的。因为所有其他被称为儿子的，是因收养的恩宠而被称为儿子，并非因本性的真实；而如果有其他被称为主的，是因被授予的权利，而非固有的。唯独基督是独一的圣子和独一的主，正如宗徒所说：\BibleQuote{只有一个主耶稣基督，万物都是藉着祂而有的。} 因此，在信经按适当顺序陈述了圣子由圣父诞生那不可言喻的奥秘之后，现在下降到祂为了人的救恩而屈尊进入的救恩计划。对于刚才称之为\Quote{天主的独生子}和\Quote{我们的主}的那一位，信经现在说：\par

9. \Quote{ \Emph{祂是由} \Emph{（藉着）} \Emph{圣神降孕，由童贞玛利亚诞生} } 这人类之中的诞生是救恩计划的方式，而之前的诞生属于天主的本体；一个是祂屈尊俯就的结果，另一个源于祂的本质本性。祂是由圣神降孕，由童贞玛利亚诞生。在此需要洁净的耳朵和纯洁的心灵。因为你们必须明白，如今在一位童贞女的子宫隐秘处，已为那一位（你们曾得知祂不可言喻地由圣父所生）建造了一座圣殿。正如在圣神的圣化中不容许有任何不完美的想法，同样在童贞女生子中也不应想象任何玷污。因为这一诞生是赐给这个世界的新生，且确实是新的。因为那位在天上是独生子的，自然地在地上也是独生子，并且是独一无二地诞生，诞生方式无与伦比，过去没有，将来也不可能有。\par

先知们关于祂的话：\BibleQuote{一位童贞女将怀孕生子}，是众所周知的，并屡次在福音中被引用。厄则克耳先知也预言了那诞生的奇妙方式，比喻性地称玛利亚为\Quote{主的门}，即主进入世界所通过的门。因为他说：\BibleQuote{那朝东的门要关闭，不可打开，没有人能通过，因为上主以色列的天主将由此通过，它将关闭。} 还有比这更明确地指向童贞女状态完好无损的保存吗？那童贞之门是关闭的；上主以色列的天主由此进入；祂由此从童贞女的子宫来到这个世界；而童贞状态被完好无损地保存，童贞女的门永远关闭。因此，圣神被说成是主肉身及其圣殿的创造者。\par

10. 由此你也可理解圣神之威严。因为福音为祂作证：当天使对童贞女说：\BibleQuote{你将生子，要给祂起名叫耶稣，因为祂要把自己的百姓从罪恶中拯救出来}时，她回答说：\BibleQuote{这事怎么成就？因为我不认识男人。}天使便对她说：\BibleQuote{圣神要临于你，至高者的能力要庇荫你，因此那要诞生的圣者，将称为天主子。}你看，天主圣三彼此合作。圣神被说成临于童贞女，至高者的能力荫庇她。这至高者的能力是什么？不就是基督自己，祂是天主的能力、天主的智慧么？这能力是谁的能力？至高者的能力。因此这里有至高者，也有至高者的能力，还有圣神。这就是天主圣三，处处隐藏又处处显现，名称和位格有别，但在天主性本体上不可分。虽然只有圣子从童贞女所生，但至高者也临在，圣神也临在，好使童贞女的受孕和生育得以圣化。\par

11. 既然这些事有先知书作为凭据，或许能堵住犹太人，尽管他们不信且怀疑。但外邦人常嘲笑我们，听我们谈论童贞女生子。因此我们必须略作答复，反驳他们的责难。我想，每一次生育都取决于三个条件：必须有成年女子，她须与男子交合，她的子宫不能荒胎。在我们所谈论的这次生育中，这三个条件缺少了一个：男子。由于我们所谈论的那位不是地上的人，而是天上的人，这欠缺由天上的圣神补足了，同时母亲的童贞得以保持不受侵犯。然而，童贞女怀孕为何被视为奇事？众所周知，东方有一种鸟称为凤凰，它无需配偶而生，或者说重生，始终保持同一，不断通过生育或重生产生自身。蜜蜂不知嫁娶，也不生产后代。还有其他一些事物也被发现遵循某种类似的生育法则。那么，由天主的能力为了更新和恢复整个世界所行的事，既然在动物的生育中能见到实例，难道反倒被认为不可信吗？然而，外邦人认为这事不可能，实在奇怪，因为他们相信自己的密涅瓦是从朱庇特的头脑中生出来的。还有什么比这更难置信、更违反自然呢？这里有一位妇女，自然的秩序得以保持，有怀孕，有适时生育；那里没有女性，只有一个男人，却——生育了！为什么相信那一个的人却对这一个感到惊奇呢？再者，他们说巴克斯神是从朱庇特的大腿生出的。这又是一个奇事，却被人相信。还有维纳斯（他们称作阿佛洛狄忒）据信是从海水的泡沫中生出，正如她的复合名称所显示。他们断言卡斯托耳和波鲁克斯是从蛋中生出，密耳弥多涅人是从蚂蚁生出。还有其他千百种事，虽反自然，却被他们信以为真，例如丢卡利翁和皮拉投掷的石头以及由此生出的庄稼人。既然他们相信如此众多的神话，那么，一位成年女子未受男子玷污，却被圣神感孕而生下神圣的后裔，这一件事难道在他们看来不可能么？如果他们难以相信，原本就不该相信那些怪奇之事，因为它们既众多又卑劣；但如果他们相信那些事，那么他们更应该接受我们的信仰，它如此尊贵和神圣，而非他们那些如此不光彩和卑陋的事。\par

12. 但他们或许会说：如果天主能使童贞女怀孕，也能使她生产，但他们认为如此伟大的存在不应以这种方式进入世界，即使没有与男子交合的玷污，分娩的过程仍会带来不体面。对此，我们简要回答，从他们的层次来回应：如果有人看见一个小孩在泥沼中将要窒息，而他自己——一个伟大有力的人——走进泥沼，几乎到了边缘，去拯救那垂死的孩子；你会指责这个人因踏进一点泥沼而被玷污呢，还是赞扬他为拯救一个垂死者的生命而怜悯呢？但所假设的是常人的情况。让我们回归到那位的本性。你认为太阳的本性比他低多少？毫无疑问，受造物比造物主低多少？现在请想：如果一束阳光照在泥沼上，它会受到任何污染吗？或者太阳因照耀污秽之物而受损吗？火也是如此，它的本性比我们所讨论的事物低多少？然而，任何物质，无论污秽或卑贱，如果接触火，都不会被认为污染火。既然物质界的情况如此明显，你能相信任何玷污和污染能够临到那超越一切火和一切光的卓越的、非物质的本质吗？最后，还要注意这一点：我们说，人是天主用地上的尘土造的。但如果天主被认为在寻求恢复祂自己的作品时被玷污，那么祂起初造物时更应被如此认为。你们无法指出祂为何做出如此违反我们羞耻感的事，却要问祂为何经历这些事，这是徒劳的。因此，使我们认为这些事如此的不是本性，而是普遍的见解。否则，身体中一切事物，既然从同一泥土形成，彼此的区别仅在于它们的用途和自然功能。\par

13. 但在解答此问题时，我们不可忽略另一项考量，即：天主的实体完全是无形体的，除非先有一种精神实体作为媒介，能够接受天主的圣神，否则它不能被引入肉体或为肉体所接受。举例来说，若我们说光能照射身体的所有肢体，但除了眼睛，没有肢体能接受光，因为只有眼睛能接受光。同样，天主圣子由童贞女所生，起初并非仅与血肉结合，而是先有一个灵魂作为肉体与天主之间的媒介被生育。因此，在灵魂作为媒介、并在理性精神的隐秘堡垒中接受天主圣言的情况下，天主由童贞女所生，并无你所想象的任何贬损。所以，在圣神圣化之处，以及那能接受天主的灵魂也参与肉体之处，不可视为有任何卑贱或不体面。在至高者的大能临在之处，不可认为任何事为不可能。在天主性的圆满充满之处，不可存有人性软弱的想法。\par

14. \Emph{他在般雀·比拉多执政时被钉十字架，并被埋葬；他下降阴府。}宗徒保禄教导我们，我们应当有\BibleQuote{我们心目的眼睛被光照}，\BibleQuote{好叫我们明了，什么是高度、宽度和深度}。\BibleQuote{高度、宽度和深度}是描述十字架：插入地中的部分他称之为深度，立在地上向上延伸的部分为高度，向左右伸展的部分为宽度。既然有如此多的死亡方式让人离世，为何宗徒希望我们心目的眼睛被光照，以明白为什么在所有死亡方式中，十字架被特别选为救主的死亡方式？那么我们必须知道，那十字架是一场凯旋，是一个得胜的旗帜。凯旋是战胜敌人的标记。当基督到来时，他立即将三个国度置于他的统治之下（这由他在\BibleQuote{致使耶稣之名，一切天上、地上和地下的事物无不屈膝}这句话中所表明），并藉他的死亡征服了这一切，于是寻找了一种与奥迹相称的死亡，使他被高举在空中，制服空中的权势，并彰显他对这些超自然和属天权势的胜利。此外，圣先知说\BibleQuote{他整日伸开双手}向着地上的人民，以便向不信者提出抗议，并邀请信者；最后，藉那插入地下的部分，他表明将阴间的国度也置于自己的统治之下。\par

15. 此外——简略触及一些更深奥的主题——当初天主创造世界时，他设立并委任了某些属天德能的权势来治理和指引必死的人类。梅瑟在\WorkTitle{申命纪}之歌中表明此事是这样进行的：\BibleQuote{至高者划分万民时，他按照天主天使的数目，划定了万民的疆界。}但其中有些，如被称为\OriginalTerm{此世的首领}{prince of this world}的，并没有按照他们领受时的律法行使天主所委托的权能，也没有教导人类遵守天主的诫命，反而教导他们追随自己邪恶的引导。因此我们被置于罪的锁链之下，因为正如先知所说：\BibleQuote{我们被卖给了我们的罪。}因为人一旦顺从私欲，就是在收取自己灵魂的代价。因此，每个人都受那些最邪恶统治者的锁链束缚；基督到来时，撕毁了这同一锁链，剥夺了他们的权势。保禄在一个伟大奥迹中表明此事，论及他说：\BibleQuote{他涂抹了那反对我们的债务文书，把它钉在十字架上，解除了率领者和掌权者的武装，公开羞辱了他们，藉着十字架向他们夸胜。}那些天主委派管理人类的统治者，变得悖逆和暴虐，开始攻击托付给他们的人，并在罪的冲突中彻底击败他们，正如厄则克耳先知神秘地暗示：\BibleQuote{在那一天，天使们要急速出来消灭埃及，在埃及的日子中他们必有惊慌；因为看，他来了。}基督剥夺了他们的大能，据说他凯旋了，并将从他们手中夺回的能力交还给人，正如他自己在福音中对门徒所说：\BibleQuote{看，我已给你们权柄，践踏蛇和蝎子，并制服仇敌的一切势力。}基督的十字架，使那些不正当地滥用所领受权威的人服从于先前被他们统治的人。至于我们，即人类，它首先教导我们要抵抗罪恶至死，并甘愿为宗教而死。其次，这同一十字架为我们树立了一个服从的榜样，正如它惩罚了那些曾是我们统治者的悖逆。请听宗徒如何藉基督的十字架教导我们服从：\BibleQuote{你们该怀有基督耶稣所怀有的心情：他虽具有天主的形体，却没有以自己与天主同等为应当把持不舍的，反而空虚自己，取了奴仆的形体，成了人的模样；既有人的样貌，就贬抑自己，服从至死，且死在十字架上。}因此，正如一位完美的师傅同时以榜样和教训来教导，基督也藉着自己首先服从至死，来教导善良的人们应付出死亡来服从的服从。\par

16. 但也许有人听到我们谈论祂的死亡会感到震惊，因为我们不久前还说祂与天主圣父是永存的，祂生于圣父的本体，并与天主圣父在主权、威严和永恒上是一体的。但不要惊慌，忠实的听者。马上你将会看到，你所听到其死亡的那一位再次成为不朽的；因为祂所屈从的死亡即将剥夺死亡。因为刚才我们所阐明的那降生成人的奥迹的目的，就是要使天主圣子的神性能力，好似隐藏在人性肉体的形状和样式下的钩子（祂正如宗徒保禄所说，\BibleQuote{成了人的形状}），引诱这世界的王来争战；向祂献上祂的肉体作为饵，其下的神性就能用钩子抓住并牢牢钩住他，藉着倾流祂无玷的血。因为唯有祂，那位不知罪污的，摧毁了所有人的罪，至少是那些用祂的血标记了信仰门框的人。因此，正如一条鱼咬住带饵的钩子，它不仅没有把饵从钩上取下，反而被拉出水面，成为他人的食物；同样，那拥有死亡权势的，在死亡中抓住了耶稣的身体，没有觉察到内藏于其中的神性的钩子，但一旦吞下它，就立刻被捉住，地狱的门闩被打破，它仿佛被从深渊中拖出，成为他人的食物。这个结果早有厄则克耳先知以同样的形象预言说：\BibleQuote{我要用我的钩子将你拉出，把你铺在地上；原野将被你充满，我要使空中所有的飞鸟都落在你身上，并且使地上所有的野兽都因你而饱饫。}达味先知也说：\BibleQuote{你打碎了那大龙的头，把它给了埃塞俄比亚的百姓作食物。}约伯也以同样的方式为此奥迹作证，因为他在主对他说话的口吻中说：\BibleQuote{你岂能用钩子将龙拉出，岂能在它的鼻子上放上嚼环？}\par

17. 因此，基督在肉身上受苦，并非由于祂的神性有所损失或贬低，而是祂的神性藉着肉身降到死亡中，为要藉肉身的软弱成就救恩；不是要按死亡的律法被死亡拘禁，而是要藉着祂的复活亲自打开死亡的门户。这就好像一位君王前往监狱，进去打开牢门，解开镣铐，打碎锁链、门闩和插销，将囚犯带出来释放，并使那些坐在黑暗中、死荫里的人重得光明和生命。因此，君王确实可以说是进了监狱，但与被关押在那里的囚犯处境不同。他们进监狱是为受罚，祂进监狱却是为释放他们免受惩罚。\par

18. 将信经传授给我们的人，以极大的远见指明了这些事情发生的时间——\Quote{在般雀比拉多治下}——以免传统因模糊和不确定而动摇。但应当知道，\Quote{祂下降阴府}这一句并未加在罗马教会的信经中，也未加在东方的教会中。不过，当说\Quote{祂被埋葬}时，似乎已经隐含了这一点。然而，以你们对\WorkTitle{圣经}的热爱和热忱，你们无疑会对我说：\Quote{这些事情应当由\WorkTitle{圣经}更明显的证据来证明。因为所当信的事情越重要，就越需要恰当而确凿的见证。}不错。但我们，因是对通晓法律的人说话，为求简洁，略去了一大片证言的森林。但如果这也要求，就让我们从许多中引用几个吧，因为我们知道，对于熟悉\WorkTitle{圣经}的人来说，一片广阔的证言之海摆在他们面前。\par

19. 首先，我们必须知道，十字架的教训并非在所有人眼中都一样。对外邦人是一回事，对犹太人又是另一回事，对基督徒又是另一回事；正如宗徒所说：\BibleQuote{我们宣讲被钉十字架的基督——对犹太人来说是绊脚石，对外邦人来说是愚妄，但对那些蒙召的，无论是犹太人还是希腊人，基督是天主的德能和天主的智慧；}又在同一处说：\BibleQuote{十字架的道理，为那些丧亡的人是愚妄，但为我们这些得救的人}，它就是\BibleQuote{天主的德能。}犹太人因从法律中得知基督应永远存在，故为祂的十字架所绊倒，因为他们不愿相信祂的复活。对外邦人来说，天主竟顺从死亡，这显得是愚妄，因为他们不明白降生成人的奥迹。但基督徒既接受了祂在肉身中的诞生与受难以及祂从死者中复活，当然相信那是战胜死亡的天主德能。\par

首先，因此，请听依撒意亚是如何预言这事的：犹太人——先知们曾向他们预言了这些事——不相信，而那些从未从先知们那里听说过的人却会相信。\BibleQuote{未曾听闻祂的，将要看见；没有听过的，将要明白。}此外，同一位依撒意亚预言：那些从幼年到老年研读法律的人不相信，而一切奥秘都应传给外邦人。他说：\Quote{万军的上主要在这座山上为万民设宴：他们要欢饮，他们要喝酒，他们要在这座山上敷油。把这些都交给万民。}这是全能者关于万民的计划。但是那些以法律知识自夸的人或许会对我们说：\Quote{你们亵渎，竟说主遭受了死亡的败坏和十字架的痛苦。}那么，请读你在\BibleBook{耶肋米亚}{哀歌}中所写的话：\BibleQuote{我们鼻中的气息、主基督，在我们的败坏中被掳；我们曾说：我们要在祂的荫下生活于外邦人中。}你听到了先知如何说主基督被掳，并且为了我们，即为了我们的罪，被交于败坏。在祂的荫下，由于犹太人持续不信，他说外邦人安息，因为我们不是生活在以色列，而是在外邦人中。\par

20. 但是，如果你们不厌倦，让我们尽可能简要地指出福音中具体的预言引文，以便那些受训于信仰基础的人将这些见证铭记在心，免得他们对自己所信之事产生任何怀疑。福音告诉我们，耶稣的朋友和同席者犹达斯出卖了祂。让我向你展示这在\BibleBook{}{圣咏}中是如何预言的：\BibleQuote{吃我面包的人，举脚踢我。}另一处：\BibleQuote{我的朋友和邻人靠近我，并反对我。}再一处：\BibleQuote{他的言语软如油，其实却是利器。}那么，他的言语软如油是什么意思？\BibleQuote{犹达斯来到耶稣跟前，对祂说：\InnerQuote{拉比，你好。}并亲吻了祂。}因此，通过亲吻的温柔谄媚，他植入了背叛的可憎利箭。主对他说：\BibleQuote{犹达斯，你用亲吻来出卖人子吗？}你注意到，祂被叛徒的贪财估价为三十块银钱。关于这点，先知也说：\BibleQuote{我对他们说：如果你们认为好，就给我工价；如果不，就算了。}紧接着：\BibleQuote{我从他们那里收了三十块银钱，将之投入上主殿中的铸炉。}这不是福音所记载的吗？犹达斯\BibleQuote{后悔自己所做的，把钱拿回来，丢在圣殿里，就出去了。}祂称其为祂的价钱，好像谴责和责备。因为祂在他们中间行了如此多的善事：使盲者看见，瘸子行走，瘫痪者行走，也使死者复生；为所有这些善事，他们用死亡作为祂的价钱，估价三十块银钱。福音中也记载祂被捆绑。这一点预言的话也通过依撒意亚预言了，说：\BibleQuote{祸哉！他们的灵魂！因为他们图谋极恶的计谋来害自己，说：让我们捆绑义人，因为他对我们无益。}\par

21. 但是有人会说：\Quote{这些事能应验在主身上吗？主能被人拘禁并拖去受审吗？}关于这点，同一位先知也会说服你。因为他说：\BibleQuote{上主亲自来到与民间的长老和首领的审判中。}那么，根据先知的见证，主受了审判，不仅受审判，还被鞭打，被人掌掴，被人吐唾沫，为我们忍受了一切侮辱和凌辱。因为所有听见宗徒们宣讲这些事的人都会极其惊讶，因此先知以他们的口吻喊道：\BibleQuote{上主，谁相信了我们的报告？}因为说天主、天主子曾受这些苦难，是难以置信的。为此，先知们预言了这些事，以免那些将要相信的人产生任何怀疑。主基督因此亲自说：\BibleQuote{我把我的背转给鞭打，我的面颊转给掌掴，我没有遮掩我的脸免受羞辱和唾沫。}此外，在祂的其他苦难中还记载，他们捆绑祂，把祂解送到比拉多那里。先知也预言了这事，说：\BibleQuote{他们捆绑了祂，把祂作为\OriginalTerm{友谊的保证}{xenium}送到雅林王那里。}但是有人反对：\Quote{比拉多不是王。}那么请听福音接着叙述的：\BibleQuote{比拉多听说祂是从加里肋亚来的，就把祂送到黑落德那里，后者当时是以色列的王。}先知正确地加上了\Quote{雅林}之名，意思是\Quote{野葡萄}，因为黑落德不是以色列家族的后裔，也不是那株主从埃及领出、\BibleQuote{栽在极肥沃的山冈上}的以色列葡萄树，而是一株野葡萄，即外来的品种。因此，他被称为\Quote{野葡萄}是完全正确的，因为他完全不是从以色列葡萄树的枝条生长出来的。先知用了\Quote{\OriginalTerm{友谊的保证}{xenium}}一词，这也符合，\BibleQuote{因为黑落德和比拉多，}如福音所证，\BibleQuote{由仇敌变成了朋友，}并且好像是他们和解的标志，双方都把捆绑的耶稣送给对方。这有何妨？只要耶稣作为救主，使不和者和解，恢复和平，并带来和谐！因此，关于这事在\BibleBook{}{约伯传}中也写着：\BibleQuote{愿主使地上君王的内心和好。}\par

22. 据说当\Person{比拉多}愿意释放祂时，全体民众喊道：\Quote{钉死祂，钉死祂！}这也由\Person{耶肋米亚}先知预言，他以主自己的口吻说：\BibleQuote{我的产业对我好像林中的狮子。它发声反对我，因此我憎恶它。所以（祂说）我抛弃并离开了我的房屋。}又在另一处说：\BibleQuote{你们向谁张口，向谁吐舌？}当祂站在审判官前时，经上记载\Quote{祂默不作声。}许多圣经为这事作证。在\BibleBook{}{圣咏}中写道：\BibleQuote{我成了一个听不见的人，嘴里没有反驳。}又说：\BibleQuote{我像一个聋子，听不见，像一个哑巴，不开口。}另一位先知又说：\BibleQuote{如同羔羊在剪毛者前，祂不开口。在祂受辱时，祂的审判被剥夺了。}记载有人给祂戴上荆棘冠冕。关于这事，在\BibleBook{}{雅歌}中听见天主圣父的声音，对\Place{耶路撒冷}加于祂儿子的侮辱感到惊异：\BibleQuote{\Place{耶路撒冷}的众女子，你们出去看看，祂的母亲给祂加冕的冠冕。}此外，另一位先知论到荆棘说：\BibleQuote{我指望它结葡萄，它却结了荆棘，且没有正义，只有哀号。}但为叫你明白这奥迹的秘密，那来除去世上罪过的一位，也必须将大地从因第一人的罪所承受的诅咒中解放出来，当时主说：\BibleQuote{因你的劳作，土地要受诅咒：它要给你生出荆棘和蒺藜。}因此，耶稣被戴上荆棘冠冕，为赦免那最初的定罪判决。祂被领到十字架，全世界的生命悬在造它的木头上。我要指出这也由先知的见证所证实。你发现\Person{耶肋米亚}如此说：\BibleQuote{来，让我们把木头放在祂的食物里，将祂从活人之地剪除。}又，\Person{梅瑟}哀悼他们说：\BibleQuote{你的生命将悬在你眼前，你日夜恐惧，不信任你的生命。}但我们必须继续，因我们已经超过了所计划的简短篇幅，以冗长的论述延长了我们的\Quote{简短话语}。然而我们还要加上几句话，免得我们似乎完全忽略了我们所承担的事。\par

23. 经上记载，当耶稣的肋旁被刺透时，\BibleQuote{祂流出了血和水。}这具有奥秘的意义。因为祂自己曾说过：\BibleQuote{从祂腹中要流出活水的江河。}祂也流出血，犹太人曾求这血归到他们和他们的儿女身上。祂流出水，为洗濯信者；祂流出血，为定不信者的罪。然而，这也可以理解为预表圣洗圣事的双重恩宠：一种是藉着水的圣洗所赐的，另一种是藉着殉道流血而求得的，因为二者都称为圣洗。但若你进一步问，为何我们的主被说成从肋旁流出血液和水，而非从其他肢体，我以为肋旁的肋骨象征了女人。既然罪与死的泉源发自第一个女人，即第一个亚当的肋骨，那么救赎与生命的泉源便从第二个亚当的肋骨涌出。\par

24. 经上记载，在我主受难时，从第六时辰到第九时辰，遍地黑暗。对此，你也会找到先知的见证：\BibleQuote{你的太阳要在中午落下。}又有\Person{匝加利亚}先知：\BibleQuote{在那一天不再有光。那一天有寒冷和冰霜，那一天为天主所认识；它既不是白天，也不是黑夜，但到了晚上，要有光。}有什么比先知的话更明白，以致他的话似乎不是对未来的预言，而是对过去的叙述？他预言了寒冷和冰霜。因为\Person{伯多禄}当时正在烤火取暖，因为天气寒冷；他感到寒冷，不仅是因为时间（早晨），也是因为他的信德。接着又说：\BibleQuote{那一天为天主所认识；它既不是白天，也不是黑夜。}什么是\Quote{既不是白天，也不是黑夜？}他难道不是明明说到白天中插入的黑暗，以及之后恢复的光明吗？那不是白天，因为它并非从日出开始；也不是完整的黑夜，因为在白天结束后，它并没有从起始得到应有的空间或延长至终点；但被恶人的罪行驱走的亮光，在晚上恢复了。因为在第九时辰之后，黑暗被驱散，太阳重新回到世界。另一位先知也为此作证：\BibleQuote{在白天，光要在天上昏暗。}\par

25. 福音接着叙述，兵士们分了耶稣的衣服，为祂的衣裳拈阄。圣神预先安排这事也由先知们预先作证，因为\Person{达味}说：\BibleQuote{他们分了我的衣服，为我的衣裳拈阄。}先知们甚至连那件紫红袍也没有沉默不语，据记载，兵士们曾戏弄地给祂穿上这件袍子。请听\Person{依撒意亚}：\BibleQuote{那从\Place{厄东}而来，身穿染红衣服的，是谁？你的衣服为何是红的，你的衣裳好像从酒榨中踏出来的一样？}祂自己回答说：\BibleQuote{我独自踏了酒榨，\Place{熙雍}女子。}因为只有祂没有犯罪，并除去了世上的罪。因为如果因一人死亡能进入世界，那么因一人——同时也是天主——生命岂不更能复得？\par

26. 又记载有人给祂饮醋，或搀了没药的酒，其苦甚于苦胆。且听先知对此的预言：\BibleQuote{他们给我苦胆作食物，我渴时，他们给我醋喝。}与此一致，梅瑟当日也对百姓说：\BibleQuote{他们的葡萄是索多玛的葡萄，他们的枝条是哥摩辣的枝条；他们的葡萄是苦胆的葡萄，他们的串是苦味的串。}再者，先知责备他们说：\BibleQuote{愚昧无知的人民，你们就这样报答主吗？}此外，在\BibleBook{}{雅歌}中也预言了同样的事，甚至指明了主被钉死的园子：\BibleQuote{我进了我的园中，我的妹妹，我的新娘，我采摘了我的没药。}先知在此清楚地描述了主所饮的搀了没药的酒。\par

27. 接下来记载：\BibleQuote{祂交付了灵魂。}这也早已由先知预言，他在圣子的位格中对圣父说：\BibleQuote{我把我的灵魂交在祢手中。}又记载祂被埋葬，一块大石头滚在墓门口。请听先知耶肋米亚对此的预言：\BibleQuote{他们将我的性命投在坑中，又用石头压在我身上。}先知的话最清楚地指向祂的埋葬。另有经文：\BibleQuote{义人被除去，免得看见邪恶；他的安息之所是平安。}在另一处：\BibleQuote{我要将恶人作为祂的坟墓；}又说：\BibleQuote{祂卧下如狮子，又如幼狮，谁能唤醒祂？}\par

28. 祂下降阴府也在\BibleBook{}{圣咏}中明确预言，其中说：\BibleQuote{祢也将我带到了死亡的尘土中。}又说：\BibleQuote{我的血有什么益处？当我下入朽坏之时？}又说：\BibleQuote{我陷入深深的淤泥，没有立足之地。}此外，若翰说：\BibleQuote{祢是那要来（无疑是指阴府）的一位，还是我们该等候另一位？}因此伯多禄也说：\BibleQuote{基督在肉身被杀，但藉着住在祂内的圣神复活了，祂降去给那些在狱中的灵魂宣讲，他们曾在诺厄时代不信从。}这里也宣告了祂在阴府所行的事。再者，主藉先知说话，如同论及将来之事：\BibleQuote{祢不会将我的灵魂撇在阴府，也不让祢的圣者见到朽坏。}同样，先知以预言的口吻论及已经实现的事：\BibleQuote{主，祢把我的灵魂从阴府领出，祢拯救了我，免我下到深渊。}接下来——\par

29. \Emph{第三日祂从死者中复活。}基督复活的荣耀使先前一切显得软弱易碎之事都大放光彩。若你先前觉得一个不朽之神不可能死亡，如今你看见，那战胜死亡并复活者不可能再是必朽的。然而你要在此理解造物主的良善：你因犯罪而堕落到何等程度，祂便降生跟随你到何等程度。不要以为全能的天主没有能力，以为祂的工程止于堕入深渊，而祂的救赎目的不能触及。我们谈论下界和上界，是因为我们被身体的固定界限所包围，被限定在我们所被指定的区域之内。但对于处处临在、无处不临在的天主，什么是下界，什么是上界呢？然而，由于取了肉身，这些也有了意义。那曾被安放在坟墓中的肉身复活了，为应验先知所说：\BibleQuote{祢不让祢的圣者见到朽坏。}因此，祂以胜利者从死者中归来，带领着阴府的掳物。因为祂领出了那些被死亡所囚禁的人，正如祂自己曾预言的：\BibleQuote{当我从地上被举起来时，我要吸引众人归向我。}福音为此作证，它说：\BibleQuote{坟墓开了，许多长眠的圣者的身体起来，显现给许多人，并进入了圣城。}无疑就是那城，宗徒所说的：\BibleQuote{那在上面的耶路撒冷是自由的，她是我们的母亲。}正如他又对希伯来人说的：\BibleQuote{原来那为万物所属、为万物所本的，要领许多儿子进入光荣，使他们救恩的元首因苦难而成为完全的，本是合宜的。}因此，祂坐在至高天上天主的右边，将那曾因第一人的过犯而死亡、却因复活之德而更新的人类肉身，经过苦难而成为完全的，安置在那里。因此宗徒也说：\BibleQuote{祂使我们一同复活，一同坐在天上。}因为祂是陶工，正如先知耶肋米亚教导的：\BibleQuote{祂亲手拿起那从祂手中坠落摔碎的器皿，又照祂眼中看为好的重新塑造。}祂看为好的是：祂所取的那必朽可坏的肉身，从石墓中复活，成为不朽不可坏的，如今不再安置在地上，而是在天上，在祂父的右边。旧约圣经充满这些奥迹。没有一位先知、一位立法者、一位圣咏作者是沉默的，几乎每一页圣经都论及它们。因此，滞留于收集证言似乎是多余的；但我们仍要引用少数几处，将那些渴望更多畅饮的人，引向神圣卷册本身的源泉。\par

30. 在\BibleBook{}{圣咏}中如此说：\newline{}\Quote{我躺下睡觉，我醒来了，因为主扶持了我。}\newline{}\newline{}又在另一处说：\newline{}\Quote{因了穷苦者的委屈和贫民的哀叹，现在我要起来，——上主说。}\newline{}\newline{}在别处，如我们以上所说：\newline{}\Quote{主啊，你把我的灵魂从地狱救出，你从下坑的人中救了我。}\newline{}\newline{}又在另一处：\newline{}\Quote{因为你转回而使我活过来，又从深渊把我提起。}\newline{}\newline{}在\BibleRef{咏 87}中，祂被最明显地论及：\newline{}\Quote{祂成了无人帮助的人，在死者中自由。}\newline{}\newline{}不是说的\Quote{一个人}，而是\Quote{如同一个人}。因为祂下降地狱时，祂是\Quote{如同一个人}；但祂是\Quote{在死者中自由}，因为死亡不能禁锢祂。因此，在同一个性体中，显示了人性软弱的能力，另一个则显示了神圣威严的能力。\newline{}\newline{}欧瑟亚先知也最明确地论及第三日：\newline{}\Quote{两天后祂要医治我们；但在第三天我们必要起来，在祂面前生活。}\newline{}\newline{}这话是祂以那些与祂一同在第三天复活、从死亡中被召回生命的人的口吻说的。他们正是那些说\Quote{在第三天我们要复活，在祂面前生活}的人。\newline{}\newline{}依撒意亚却明白地说：\newline{}\Quote{祂从地上领出了羊群的大牧者。}\newline{}\newline{}此外，妇女们要看见祂的复活，而经师和法利塞人及百姓却不信，这也是依撒意亚预言的话：\newline{}\Quote{你们这些从观看中出来的女子，来吧！因为这是一群没有智慧的百姓。}\newline{}\newline{}至于那些妇女，记载说她们在复活后去了坟墓，寻找祂却没有找到，如同玛利亚玛达肋纳，记载说她在天亮之前来到坟墓，没有找到祂，哭着对那里的天使说：\Quote{他们把主搬走了，我不知道他们把祂放在哪里了}——连这事也在\BibleBook{}{雅歌}中预表了：\newline{}\Quote{我在床上寻觅我心爱的；我在夜间寻觅祂，却没有找着。}\newline{}\newline{}至于那些找到祂并抱住祂脚的人，在同一书中预言说：\newline{}\Quote{我要抱住我心爱的，不放祂走。}\newline{}\newline{}这些经文是从许多经文中选取的少数；因我们追求简练，不能堆积更多。\par

31.\newline{}\Emph{祂升了天，坐在圣父的右边：从那里祂要来审判生者死者}。\newline{}\newline{}这些条文以适当的简练方式出现在信经中论述圣子的部分之末。所说的是清楚的，但问题是如何和在何种意义上理解。因为\Quote{升}、\Quote{坐}和\Quote{来}，除非你按照天主性尊严来理解这些词，否则似乎指向某种人性的软弱。因为在完成了地上应做之事，并把灵魂从地狱的囚禁中召回之后，祂被说成升到天上，正如先知所预言：\Quote{祂上升高处，掳掠了俘虏，赐予人恩惠}——即那些恩惠，就是伯多禄在宗徒大事录中关于圣神所说的：\Quote{祂既被举扬到天主的右边，就领受了所恩许的圣神，并倾注了你们所见所闻的恩惠。}祂将圣神的恩惠赐予人，因为基督藉着从死者中复活，将以前被魔鬼藉罪带入地狱的俘虏召回天上。\newline{}\newline{}因此祂升了天，不是到天主圣言以前不在的地方，因为祂永远在天上，且住在父内，而是到了圣言成血肉以前未曾坐过的地方。最后，因为进入天门对天朝的臣仆和王子来说似乎是新事，他们看见血肉之躯穿透天国的隐秘内室，便彼此说，正如充满圣神的达味所宣告的：\Quote{城门，请高举你们的门楣，古老的门户，请高起，光荣的君王要进来。谁是这位光荣的君王？是英勇有力的上主，是战争中的勇者。}这些话不是指着天主性的能力说的，而是指着血肉上升至天主右边的新奇事。同一位达味在别处说：\Quote{天主在欢呼声中上升，上主在号角声中上升。}因为胜利者惯于在号角声中从战场归来。关于祂，也说道：\Quote{祂在天上建造自己的上升。}又：\Quote{祂上升于革鲁宾之上，驾着风的翅膀飞行。}\par

32. 坐在圣父的右边是降生成人奥迹的一部分。因为那无形体的本性，在没有取得肉身的情况下，是不适宜的；天上的座位之尊荣，不是为求于天主性，而是求于人性。因此经上论及祂说：\Quote{天主，祢的宝座从那时起就已预备；祢从永远就有。} 所以，主耶稣将要坐的宝座，从永远就已预备好了，\Quote{因耶稣之名，一切在天上的、地上的和地底下的，无不屈膝叩拜；众口皆明认耶稣基督是主，以光荣天主圣父。} 达味也这样论及祂说：\Quote{上主对我主说：你坐在我右边，等我使你的仇敌作你的脚凳。} 主在福音中论及这些话，对法利塞人说：\Quote{那么，达味在圣神内怎会称祂为主，祂又怎能是达味之子？} 由此祂显示：按圣神，祂是主，按肉身，祂是达味之子。因此主又在另一处说：\Quote{我实在告诉你们：从今以后，你们要看见人子坐在天主大能的右边。} 宗徒伯多禄论及基督说：\Quote{祂已升天，坐在天主的右边。} 保禄写信给厄弗所人也说：\Quote{祂已使基督从死者中复活，叫祂坐在自己的右边。}\par

33. 祂将要审判活人死人，我们由许多圣经的见证得知。但在我们引用先知们关于这点的言论之前，我们认为有必要提醒你们：这信德的道理要我们每日对审判者的来临保持警醒，好使我们如此行事，如同要向那临近的审判者交账一样。这正是先知论及那有福之人所说的：\Quote{他在审判中安排自己的言语。} 然而，当祂被说成审判活人死人时，这并不意味着有些人还活着而另一些人已经死了，有些人将来接受审判；而是祂将审判灵魂和身体两者，其中灵魂意指\Quote{活人}，身体意指\Quote{死人}；正如主自己在福音中所说：\Quote{你们不要害怕那些能杀死身体，却不能杀害灵魂的；却要害怕那能使灵魂和身体都丧亡在地狱中的。}\par

34. 现在，我愿简要向你表明，这些事早已由先知们预言了。你若愿意，自可从更丰富的圣经中收集更多。\Person{玛拉基亚}先知说：\Quote{看，全能的上主将来到，谁能忍受他来临的日子，谁能承受看见他的情景？因为他来临有如炼金之人的火，又像漂布之人的碱。他必坐下，如同炼净金银的人。} 但为使你更确切地知道所说的是哪一位主，请听\Person{达尼尔}先知也预言说：\Quote{我看见，}他说，\Quote{在夜间的异象中，看，有一位像\OriginalTerm{人子}{Son of Man}的，乘着天上的云而来，他走近\OriginalTerm{万古常存者}{Ancient of days}，被引到他面前；他便赐给他统治权、尊荣和国度。一切民族、部族、语言都事奉他。他的统治是永远的统治，不消逝；他的国度永不灭亡。} 这些话不仅教导我们他的来临和审判，也教导他的统治和国度：他的统治是永恒的，他的国度是不可摧毁的，没有终结；正如信经所说：\Quote{他的国度将没有终结。} 因此，若有人说基督的国度终有一日会结束，那就离信德甚远。然而，我们应当知道，仇敌惯常以狡猾的欺诈，假冒基督这救恩的来临，以欺骗信徒，并代替那在父的威严中期待来临的\OriginalTerm{人子}{Son of Man}，准备那\OriginalTerm{丧亡之子}{Son of Perdition}，行奇迹和虚假征兆，以便将\OriginalTerm{假基督}{Antichrist}引入世界，取代基督；主自己在福音中预先警告了犹太人：\Quote{我因我父的名而来，你们却不接待我；另有一个会因自己的名而来，你们反要接待他。} 又说：\Quote{当你们看见\Person{达尼尔}先知所说的\OriginalTerm{招致荒凉的可憎之物}{abomination of desolation}，立于圣地时，读者须要明白。} 因此，\Person{达尼尔}在他的异象中非常详尽地述说了那\OriginalTerm{欺骗}{delusion}的来临；但没有必要再举例，因为我们已足够扩展；因此，任何愿意了解更多这些事的人，可自行查阅那些异象。宗徒自己也说：\Quote{不要让任何人以任何方式欺骗你们，因为那日子不会来，除非先有背叛之事，且那\OriginalTerm{罪恶的人}{man of sin}被启示出来，那\OriginalTerm{丧亡之子}{Son of Perdition}，他反对并高举自己超过一切称为神或受人敬拜的，以致坐在天主的殿中，显示自己好像是神。} 不久之后：\Quote{那时那不法者将被启示出来，主耶稣要以口中的气息杀死他，并以来临的荣光消灭他：他的来临是出于撒殚的运作，具有一切权能、征兆和欺骗的奇迹。} 又过不久：\Quote{因此，主要把强烈的\OriginalTerm{欺骗}{delusion}送到他们中间，使他们相信谎言，使所有不信真理的人都受审判。} 因此，这\Quote{\OriginalTerm{欺骗}{delusion}}由先知、福音作者和宗徒们的话预先告知我们，以免有人将\OriginalTerm{假基督}{Antichrist}的来临误认为基督的来临。但正如主自己所说：\Quote{当他们对你们说：看，基督在这里，或看，他在那里，你们不要相信。因为许多假基督和假先知将要起来，欺骗许多人。} 然而，让我们看他如何指明真基督的审判：\Quote{如同闪电从东方闪到西方，\OriginalTerm{人子}{Son of Man}的来临也将如此。} 因此，当真主耶稣基督来临时，他必坐下设立审判的宝座。如他在福音中所说：\Quote{他要把绵羊和山羊分开，}即义人与不义之人；正如宗徒所写：\Quote{我们都必须站在\OriginalTerm{基督的审判台}{judgment-seat of Christ}前，好使每人按他在身体上所行的，或善或恶，领取应得的赏报。} 此外，审判不仅针对行为，也针对思想，如同一位宗徒所说：\Quote{他们的思想互相控告或辩护，在天主审判人隐秘事的日子。} 但关于这些，说这些就够了。信德顺序接着是——\par

35. \Emph{及在圣神内}。之前关于基督较为详尽的论述，涉及他降生成人及苦难的奥迹，并且由于这样插入，属于他的位格，而稍微延迟了提及圣神。否则，如果只考虑天主性，如同在信经开头我们所说\Quote{我信全能的天主圣父}，随后\Quote{在耶稣基督、他的独生子、我们的主内}，同样我们加上\Quote{及在圣神内}。但上述关于基督的一切细节，如我们所说，都与\OriginalTerm{肉身的安排}{dispensation of the flesh}（即他的降生成人）有关。借着提及圣神，\OriginalTerm{天主圣三}{Trinity}的奥迹得以完成。因为既然提到一位圣父，且没有别的圣父；一位\OriginalTerm{独生子}{only-begotten Son}被提到，且没有别的独生子；同样，有一位\OriginalTerm{圣神}{Holy Ghost}，且不能有另一位圣神。因此，为使\OriginalTerm{位格}{Persons}可以区分，表达关系的术语（属性）各有不同，由此第一位被理解为圣父，万物来自他，他本身也没有父；第二位是圣子，由圣父所生；第三位是圣神，\OriginalTerm{由父和子所共发}{proceeding from both}，并圣化一切。但为在圣三中显示同一个天主性，既然冠以介词\Quote{在}，我们说我们相信\Quote{\Emph{在}天主圣父内}，同样我们也说\Quote{\Emph{在}基督、他的圣子内}，同样\Quote{\Emph{在}圣神内}。但我们的意思在后续中将更加清楚。因为信经接着是说——\par

36. \Quote{\Emph{圣而公教会；罪的赦免，肉身的复活}}。经文不说：\Quote{信圣而公教会}，也不说：\Quote{信罪的赦免}，也不说：\Quote{信肉身的复活}。因为若加了\Quote{信}这一介词，它就会与前几款具有同样的效力。但在如今这些宣示关于天主性的信仰的条款中，我们却说：\Quote{信天主圣父}、\Quote{信耶稣基督祂的圣子}和\Quote{信圣神}；而在其余涉及非天主性、而是受造物和奥迹的条款中，则不加\Quote{信}这一介词。我们不说：\Quote{我们信圣而公教会}，而是说：\Quote{我们信圣而公教会}，不是作为天主，而是作为聚集归于天主的教会；我们相信有\Quote{罪的赦免}，不说：\Quote{我们信罪的赦免}；我们相信将有\Quote{肉身的复活}，不说：\Quote{我们信肉身的复活}。因此，正是通过这个单音节介词，造物主与受造物得以区分，属神的事与属人的事得以分离。\par

这圣神在旧约中默感了法律和先知，在新约中默感了福音和宗徒书信。因此宗徒也说：\Quote{\BibleQuote{凡受天主默感所写的圣经，有益于教诲}}。故此，在此处依从教父的传统，列出新旧约中那些根据先辈传承、被认为是受圣神默感并交付给基督的教会的经书，似乎恰如其分。\par

37. 因此，旧约首先传下来的是梅瑟五书：\BibleBook{}{创世纪}、\BibleBook{}{出谷纪}、\BibleBook{}{肋未纪}、\BibleBook{}{户籍纪}、\BibleBook{}{申命纪}；接着是\BibleBook{}{若苏厄书}（努恩之子若苏厄）、\BibleBook{}{民长纪}连同\BibleBook{}{卢德传}；然后是四卷\BibleBook{}{列王纪}（统治纪），希伯来人将其算作两卷；遗漏之书，即\BibleBook{}{编年纪}（日纪），以及两卷\BibleBook{}{厄斯德拉书}（\BibleBook{}{厄斯德拉}和\BibleBook{}{乃赫米雅}），希伯来人将其算作一卷，还有\BibleBook{}{艾斯德尔传}；先知书方面：\BibleBook{}{依撒意亚}、\BibleBook{}{耶肋米亚}、\BibleBook{}{厄则克耳}和\BibleBook{}{达尼尔}；此外，十二小先知合为一卷；\BibleBook{}{约伯传}和\BibleBook{}{达味圣咏}各为一卷。\Person{撒罗满}给予教会三卷书：\BibleBook{}{箴言}、\BibleBook{}{训道篇}、\BibleBook{}{雅歌}。这些涵盖了旧约的经书。\par

新约则有四福音：\BibleBook{玛窦}{福音}、\BibleBook{马尔谷}{福音}、\BibleBook{路加}{福音}、\BibleBook{若望}{福音}；\BibleBook{路加}{宗徒大事录}，由\Person{路加}撰写；\Person{保禄}宗徒的十四封书信、\Person{伯多禄}宗徒的两封信、\Person{雅各伯}——主的兄弟和宗徒——的一封信、\Person{犹达}的一封信、\Person{若望}的三封信、\BibleBook{若望}{默示录}。这些是教父们纳入正典的经书，并愿我们从中推导出我们信仰的证据。\par

38. 但应当知道，还有其他一些经书，教父们不称之为\OriginalTerm{正典}{Canonical}，而称为\OriginalTerm{教会书}{Ecclesiastical}：即\BibleBook{}{智慧篇}，称为\WorkTitle{撒罗满的智慧}；另一\BibleBook{}{智慧篇}，称为\WorkTitle{西拉赫之子的智慧}，拉丁文统称为\BibleBook{}{德训篇}，这名称并非指书的作者，而是指写作的性质。同属这一类的还有\BibleBook{}{多俾亚传}、\BibleBook{}{友弟德传}、\BibleBook{}{玛加伯书}。在新约中，有一本小书称为\WorkTitle{赫尔马牧者书}，[以及]所谓的\WorkTitle{两条路}或\WorkTitle{伯多禄的审判}；所有这些书，教父们愿意在教会中诵读，但不可用于教义的确认。其他著作则称为\OriginalTerm{伪经}{Apocrypha}，他们不愿在教会中诵读。\par

这些是教父们传给我们的传统，如前所述，我认为在此处阐明它们时机恰当，为的是教导那些正在学习教会和信仰基本知识的人，使他们知道应从哪些天主圣言的泉源汲取活水。\par

39. 按照信仰的次序，接下来我们论及\Emph{圣而公教会}。我们已在前面提及，为何信经在此处不像前一款那样说\Quote{信圣而公教会}。因此，那些在上文受教，在圣三奥迹中信一天主的人，也必须相信这一点：有一个圣而公教会，其中有同一个信仰、同一个圣洗，其中信一天主圣父、一主耶稣基督祂的圣子、一圣神。这圣而公教会是毫无玷污或皱纹的。因为有许多其他人聚集教会，如\Person{马西翁}、\Person{瓦伦提努}、\Person{厄俾翁}、\Person{摩尼}和\Person{亚略}，以及所有其他异端者。但那些教会并非毫无不信的玷污或皱纹。因此先知论及他们说：\Quote{\BibleQuote{我恼恨好党的集会，也决不与恶人同席}}。但论到这保存基督完整信仰的教会，请听圣神在\BibleBook{}{雅歌}中说的：\Quote{\BibleQuote{我的鸽子，我的完人，只有一个}}。那么，凡在教会中领受这信仰的人，不要让他在虚妄的集会中转向，也不要与那些行不义的人同流。\par

因为马西昂的集会是一个\OriginalTerm{虚妄的集会}{Council of vanity}，他否认基督的父是天主、造物主，祂藉自己的圣子创造了世界。厄比雍的集会是一个虚妄的集会，因为他教导说，我们信基督的同时，还必须遵守肉身的割损礼、安息日的停工、惯常的祭献以及法律字面上的所有其他规条。摩尼的集会是一个虚妄的集会，就其教导而言：首先他称自己为\OriginalTerm{护慰者}{Paraclete}，其次他说世界是由一个邪恶的神创造的，否认天主造物主，拒绝旧约，主张两种本性，一善一恶，彼此对立，断言人的灵魂与天主同永恒，按照毕达哥拉斯派学说，它们通过各种出生的循环进入牛、动物和野兽，否认我们肉身的复活，坚持主的受难和诞生不是在肉身的真实中，而只是表象。当撒摩撒塔的保禄和他的继任者福提诺后来教导说，基督不是由圣父在万世之前所生，而是从玛利亚开始的，他们不信祂是神而成为人，而是从人成为神，这是虚妄的集会。当亚略和欧诺弥教导说，天主子不是从圣父的实体本身所生，而是从无中受造，并且天主子有一个开始，且低于圣父，这是虚妄的集会；此外他们还断言圣神不仅低于圣子，而且是一个服役之神。那些承认圣子确实出于圣父的实体，却区分并分离圣神的人，他们的集会也是虚妄的，然而救主在福音中显示圣三的能力和神性是同一的，说：\BibleQuote{你们要去使万民成为门徒，因父、及子、及圣神之名给他们授洗}，而人将天主所结合的分离，显然是亵渎。那由顽固恶毒的争辩从前聚集起来的集会也是虚妄的，它肯定基督确实取了人的肉身，却没有同时取理性的灵魂，因为基督将同一救恩赐予了肉身、动物的灵魂以及人的理智和心灵。那多纳图在整个非洲召集的集会也是虚妄的，他以\OriginalTerm{交付圣书者}{traditorship}的罪名控告教会；诺瓦都也用这集会扰乱人心，他否认对跌倒者赐予悔改，并谴责第二次婚姻，尽管可能是不得已而缔结的。所有这些，都要躲避，如同恶人的集会。还有那些人，若有的话，据说断言天主子看不见也不认识圣父，如同祂自己被圣父认识看见一样；或者基督的王国将有终结；或者肉身不会在其本体的完全恢复中复活；那些人还否认天主对所有人有公义的审判，并断言魔鬼将被免除其应受的永罚之刑。对所有这些，我再说，让信徒充耳不闻。但要持守圣教会，她承认全能的天主圣父、祂的独生子耶稣基督我们的主、以及圣神，具有同一和谐一致的本体，相信天主子由童贞女所生，为人的救恩而受苦，在祂所生的同一肉身中从死者中复活；最后，盼望祂将作为众人的审判者降临，藉着祂，也宣讲了\Emph{罪过的赦免与肉身的复活}。\par

40. 至于\Emph{罪过的赦免}，简单的相信就足够了。因为当一位君王施恩时，谁会问其原因或理由呢？当一位人间君主的宽厚不宜被讨论时，人的鲁莽岂可讨论天主的慷慨？因为异教徒常嘲笑我们，说我们自欺，以为行为所犯的罪行可以藉言语洗净。他们说：\Quote{难道杀人者可以不再是杀人者，犯奸者可以算不是犯奸者吗？那么，犯这类罪的人怎能一下子成为圣洁呢？}但对此，如我所说，我们藉信德而非理性来回答更好。因为应许此事的是万王之王，保证此事的是天地的主宰。你难道要我拒绝相信，那位用地上的尘土造我为人，能使我有罪的人成为无罪？那位当我瞎时使我看见，聋时使我听见，瘸时使我行走，能为我恢复失去的纯洁吗？再诉诸自然的证言——杀人并非总是有罪的，但出于恶意、不依法而杀人，那是有罪的。因此，在这些事上，定我罪的并非行为，因为有时行为是正确的，而是心灵的恶意。如果我的心灵曾变得有罪，罪在其中起源，如今被改正了，那么为什么在你看来，我以前有罪，现在就不能成为无罪呢？因为如果我已表明，罪行不在于行为而在于意志，正如一个邪恶的意志，受到恶神的驱使，使我有罪并该死，那么，受到善神的驱使，意志转为良善，便使我恢复纯洁与生命。在其他所有罪行中也是如此。这样，我们的信德与自然理性之间毫无冲突，因为罪过的赦免不归于行为——行为一旦作出便不能改变——而归于心灵，可以肯定，心灵能从恶转向善。\par

41. 这最后一条，肯定\Emph{肉身的复活}，以简洁的言辞总结了所有圆满的总和。虽然在这一端上，教会的信仰也受到攻击，不仅来自外邦人，也来自异端者。因为瓦伦提努斯完全否认肉身的复活，摩尼教徒也是如此，正如我们上面所表明的。但他们拒绝听从依撒意亚先知的话，他说：\BibleQuote{死去的人要复活，在坟墓里的人要兴起。}或最明智的达尼尔，他宣告：\BibleQuote{那时，许多睡在尘土中的人要醒来，这些要进入永生，而那些要蒙受羞辱和混乱。}然而，即使在福音书中（他们似乎接受福音），他们也应该从我们的主和救主那里学到，当祂教训撒杜塞人时说：\BibleQuote{关于死人的复活：你们没有读过天主在荆棘中对你们所说的吗？我是亚巴郎的天主，依撒格的天主，雅各伯的天主。他不是死人的天主，而是活人的天主。}在那里，在前文中祂宣告了复活的光荣是何等伟大，说：\BibleQuote{但在死人的复活中，他们不娶也不嫁，而是如同天上的天使一样。}但复活的功能赋予人天使般的状态，使他们从地上复活后不再与牲畜一起生活在地上，而是与天使一起在天上——只是那些更纯洁的生活使他们适合于此的人——即那些即使在现在，也以贞洁保全他们灵魂的肉身，并使它服从于圣神，从而借着圣化的能力，除去了所有罪的污秽，转化为属灵的光荣，有资格被接纳进入天使的团体。\par

42. 但不信的人喊叫说：\Quote{那已经腐烂、分解、变为尘土、有时被大海吞没、被波浪冲散的肉，怎能再被聚集，再次成为一体，人身体从其中重新形成呢？}对这些人，我们的第一个回答是保禄的话：\BibleQuote{愚昧的人啊！你所播种的，若不先死，就不会活过来。你所播种的，不是那将来要长成的形体，只是赤裸的谷粒，或是麦子，或是别的种子；但天主随自己的心意给它一个形体。}难道你不相信，你每年看见发生在你撒在地里的种子中的事，也要发生在你的肉身上（这肉身按天主的法律被种在地里）吗？请问，你为什么对天主的能力有如此低下的看法，以致不相信可能将构成每个人肉身的散落尘土重新收集并恢复其原来的结构？当你看到凡人的智慧寻找深埋地下的金属矿脉，有经验的眼睛发现黄金，而无经验的人以为只是泥土时，你难道拒绝承认这一事实吗？我们为什么拒绝将这些事归于创造人的天主，而祂所造的人却能做到这么多？当凡人的智慧发现黄金有其特有的矿脉，银有另一种，铜有远为不同的矿脉，铁和铅有各异且分明的矿脉隐藏在那看似泥土的表层之下时，难道神圣的能力被认为无法发现和区分属于每个人肉身的组成微粒，即使它们似乎已经分散了吗？\par

43. 但让我们尽力通过自然的理由帮助那些信仰软弱的灵魂。假如有人将不同种类的种子混合在一起，随意撒在地里，那么每一种谷粒，无论撒在哪里，不都会在适当的时候，按照其自身的特定本性发芽，以再现其自身形状和自身身体的状态吗？\par

因此，每个个别肉身的实体，虽然其微粒已经以各种不同的方式分散，但其中有一个不朽的原则，因为它是永生灵魂的肉身，而在天主按其善意所定的时间，它自己的组成微粒将从地上被收集并吸引过来，将恢复成死亡先前所解散的那个形状。这样，每个灵魂将恢复的，不是一个混乱的或异己的身体，而是它活着时所拥有的自己的身体，为了使肉身连同自己的灵魂，为现世生命的斗争，要么因无玷而受冠冕，要么因污秽而受惩罚。因此，我们的教会，在传授信仰时，不用其他教会所传的\WorkTitle{信经}中\Quote{肉身的复活}一词，而是谨慎地加上了代词\Quote{这个}——\Quote{\Emph{这个}肉身的复活}。\Quote{这个}无疑是指那宣认\WorkTitle{信经}的人，他一边念这个词，一边在额头上画十字圣号，使每个信徒知道：他的肉身，若在罪中保持洁净，将成为荣誉的器皿，为主有用，预备行各样的善事；但若被罪污秽，则将成为义怒的器皿，注定毁灭。\par

但现在，关于复活的荣耀以及天主以自身所许诺的伟大，若有人愿意更详尽地知晓，他几乎能在所有神圣经卷中找到相关论述；我们从这些经卷中，仅以提醒的方式，在此处提及几段，然后就结束你所托付的这部作品。宗徒保禄在断言会死的肉身将要复活时，使用了如下的论证：\BibleQuote{若死者不复活，基督也就没有复活。若基督没有复活，我们的宣讲便是空的，你们的信德也是空的。}紧接着：\BibleQuote{但是现今基督已从死者中复活，成为安眠者的初果。因为死亡既由一人而来，死者的复活也由一人而来。就如所有的人都因亚当而死，同样所有的人都因基督而活。但各人要按自己的次序：基督是初果，以后是那些属于基督的，在祂来临时；再后才是结局。}随后他又说：\BibleQuote{看，我告诉你们一个奥秘：我们固然都要复活，但不全都改变；}或按其他抄本：\BibleQuote{我们固然都要安眠，但不全都改变；在一刹那，一眨眼之间，在最后的号角声中；因为号角一响，死者将以不朽坏的身体复活，我们也要改变。}然而，无论哪一个才是正确的经文，他在给得撒洛尼人的书信中写道：\BibleQuote{弟兄们，关于安眠的人，我不愿你们不知道，以免你们忧伤，像那些没有希望的人一样。因为我们若信耶稣死了，也复活了，同样，那些通过耶稣安眠的人，天主也必把他们与耶稣一同带来。我们凭主的话告诉你们这一点：我们这些活着存留到主来临时的人，绝不会在安眠的人之先。因为主必亲自从天降下，有呼喊声，有总领天使的声音，并有天主的号角；那些在基督内死了的人要先复活；然后我们这些活着存留的人，将同他们一起被提到云中，在空中迎接基督，这样我们就常与主在一起。}\par

44. 但为使你不以为这是保禄独有的新教义，我也要引述厄则克耳先知藉圣神所预言的话。他说：\BibleQuote{看，我要打开你们的坟墓，把你们从你们的坟墓里领出来。}再让我回忆约伯，他充满奥秘言语，是如何明白地预言死者的复活。\BibleQuote{树木还有希望；因为若被砍下，它还会发芽，它的嫩枝决不会止息。但若它的根在土中衰老，它的树干在尘土中死去，然而一闻到水气，它就会重新发旺，如同新栽的植物一样。但人若死了，他就离去了吗？会死的人若倒下了，他就不再有了吗？}你难道没有看见，在这些话中，他仿佛是在诉诸人的羞耻感，说：\Quote{人类竟如此愚蠢，当他们看到被砍下的树桩又从地中发芽，枯木又重新复活，却以为自己的情况与树木全无相似之处吗？}但为使你相信约伯的话应作为问句来读，当他说：\BibleQuote{但会死的人倒下了，难道他不再起来吗？}请从下文取这个明证；因为他随即加上：\BibleQuote{但人若死了，他还能活吗？}紧接着他又说：\BibleQuote{我要等候，直到我重新被造；}随后他又重复同样的话：\BibleQuote{那如今正承受这苦杯的我的皮肤，谁能在世上使它再起来呢？}\par

45. 以上足够证明我们在信经中所宣认的\Quote{这个肉身的复活}。至于添加\Quote{这个}一词，看它与我们从圣经中所引证的一切是多么一致。约伯在他上面所解释的经文\BibleQuote{祂将再举起我的皮肤，就是如今正承受这苦杯的}——即正在经受这些苦难的——又指明了别的什么呢？他岂不是明说这个肉身的复活，即现在正经历极端考验和患难的这肉身吗？此外，当宗徒说：\BibleQuote{这必朽坏的必须穿上不朽坏的，这必死的必须穿上不死的}，他的话岂不像是触摸自己的身体、用手指着它说的吗？因此，这个现在必朽坏的身体，将因复活的恩宠成为不朽坏的；这个现在必死的，将穿上不死的德能，以使\BibleQuote{基督从死者中复活，不再死，死亡再不能管辖祂}；同样，那些在基督内复活的人，将永不再感到腐朽或死亡，并非因为肉身的本性被抛弃，而是因为其状态和品质已被改变。因此，将有一个身体从死者中复活，是不朽坏、不死的，不仅对义人如此，对罪人也如此；对义人，使他们能永远与基督同在；对罪人，使他们无终地承受应受的惩罚。\par

46. 义人将永远与我们的主基督同在，这一点我们已在上面证明，在那里我们指出，宗徒说：\Quote{然后我们这些活着存留的人，将同他们一起被提到云中，在空中与基督相遇，这样我们就常与主同在。} 不要惊奇，圣人们的肉躯将在复活时被改变成如此光荣的状态，以致被提到空中与天主相遇，悬在云中，飘在空中，因为同一位宗徒阐述天主赐予爱祂之人的伟大事物时说：\Quote{祂将改变我们卑贱的身体，使之与祂的光荣身体相似。} 因此，若说圣人们的身体被升到空中，这绝无荒谬之处，因为他们被说成是依照基督身体的肖像而更新，而基督的身体坐在天主的右边。但圣宗徒还补充了这一点，不论是指他自己，还是指与他同地位或同功勋的其他人：\Quote{祂将使我们与基督一同复活，一同坐在天上。} 因此，既然天主的圣人们拥有这些关于义人复活的应许以及无数类似的应许，那么现在相信先知们所预言的也就不难了，即：\Quote{义人将在天主的国里发光如同太阳，如同穹苍的光辉。} 因为，谁会觉得这困难呢？那些将来在天上被预备享有天主天使的生命和交往的人，或者那些被描述为将来要相似于基督身体荣耀的人，他们怎么能没有太阳的光辉，不能以星辰和穹苍的灿烂为装饰呢？关于这荣耀，救主亲口应许了，圣宗徒说：\Quote{播种的是属魂的身体，复活的是属神的身体。} 因为如果这是真的——它确实是真的——天主将赏赐每一位义人和圣人与天使为伴，那么祂也必定会将他们的身体改变成属神身体的荣耀。\par

47. 不要让你觉得这个应许与身体的自然构造相悖。因为如果按经上所载，我们相信天主取地上的泥土造了人，我们身体的起源是天主用泥土变化成了肉，那么，基于同样的原理，既然泥土被说成是提升到属魂的身体，那么属魂的身体反过来又被相信是提升到属神的身体，这为什么在你看来是荒谬或不合情理的呢？这些事以及许多类似的事，你会在神圣的圣经中找到，关于义人的复活。复活时，罪人也将得到不朽和不死的状态，正如我们上面所说的，为的是天主将这状态赐给义人作为永远的荣耀，而赐给罪人作为羞耻和惩罚的延长。因为上文我们引用过的先知的话也清楚地说：\Quote{许多人将从地上的尘土中起来，有的进入永生，有的进入羞耻和永远的耻辱。}\par

48. 如果我们已理解，全能的天主以何等庄严的意义被称为圣父，我们的主耶稣基督以何等奥秘的含义被认为是祂的独生子，祂的圣神以何等圆满的意义被称为圣神，以及天主圣三如何是同一本体却有关系与位格的区分；还有，由童贞女的诞生是什么，圣言在肉身上的诞生是什么，十字架的奥秘是什么，我们主下降阴府的目的何在，复活的荣耀是什么，灵魂从阴间的囚禁中被解救是什么，祂的升天是什么，审判者的来临是什么；此外，如何应承认圣教会与虚妄的集会相对立，圣卷的数目是多少，异端的集会应当避免，在罪的赦免中，天主的自由与人的自然理性之间并无任何对立，以及如何不仅神圣的神谕，而且主和救主自身的榜样以及自然理性的结论，都确认我们肉身复活的真理——如果我说，我们按照前面阐述的传统规则明智地依次理解了这些，那么我们恳求主赐予我们和所有听到这些话的人，使我们保持所领受的信德，完成我们的赛程之后，能等待为我们存留的正义之冠，并被算在那些复活进入永生的人中，从羞耻和永远的耻辱中被拯救出来，藉着我们的主基督，藉着祂，荣耀与权能归于全能的天主圣父及圣神，至于永远。阿们。\par

\SourceRecord{Translated by W.H. Fremantle.；Nicene and Post-Nicene Fathers, Second Series；Vol. 3.；Edited by Philip Schaff and Henry Wace.；Buffalo, NY: Christian Literature Publishing Co.,；1892.；<http://www.newadvent.org/fathers/2711.htm>.}
